\documentclass[10pt,a4paper]{amsart}
\usepackage[margin=19mm]{geometry}
\usepackage{amsmath,amssymb,mathtools}
\usepackage[scr=rsfs,cal=euler]{mathalfa}
\usepackage{xfrac}
\usepackage[unicode,hidelinks,bookmarks=false]{hyperref}
\newcommand{\ie}{i.e.\ }
\newcommand{\cf}{cf.\ }
\newcommand{\resp}{respectively\ }
\newcommand{\Subsection}{\subsection}
\providecommand{\nobreakdash}{\nobreak\hbox{-}}
\setlength{\emergencystretch}{2em}
\multlinegap0pt
\usepackage{amssymb}
\usepackage{xcolor}
\usepackage{algorithm}\usepackage{algpseudocode}
\usepackage{mdframed}
\usepackage{tikz}
\newtheorem{lemma}{Lemma}
\newtheorem{proposition}{Proposition}
\newcommand{\Ber}{Ber}
\newcommand{\F}{\mathcal{F}}
\newcommand{\KL}{\mathrm{KL}}
\newcommand{\Pc}{\mathcal{P}}
\newcommand{\TV}{\mathrm{TV}}
\newcommand{\Unif}{\mathrm{Unif}}
\newcommand{\X}{\mathcal{X}}
\newcommand{\Y}{\mathcal{Y}}
\newcommand{\eps}{\varepsilon}
\title{Uniform Laws of Large Numbers in Product Spaces}
\author{Ron Holzman, Shay Moran, Alexander Shlimovich}
\date{Source: arXiv:2603.24493v1 (CC BY 4.0)}
\begin{document}
\maketitle

\noindent\emph{Source and licence.} Uniform Laws of Large Numbers in Product Spaces. Version arXiv:2603.24493v1 (CC BY 4.0), distributed by arXiv under the Creative Commons Attribution 4.0 International licence, http://creativecommons.org/licenses/by/4.0/, as recorded in the arXiv licence metadata for this exact version and shown on the abstract page https://arxiv.org/abs/2603.24493v1. Full licence notice: \texttt{licenses/arxiv-2603.24493-CC-BY-4.0.txt}.\par\smallskip

\noindent\textit{Editorial scope.} Counted unit: the article's proposition prop:Omega-d-product (source0.tex lines 3635--3665). The article's complete original proof of it is delivered with it (source0.tex lines 3667--3828, one \texttt{proof} environment of 162 lines). No mathematics is written, repaired or re-derived: the statement and the proof are the article's own lines, in the article's own notation, and no retained line is substituted or re-flowed.  Retained uncounted as the source-local context the proof needs: lines 3400--3425; lines 3462--3489; lines 3558--3581.  Nothing was omitted from any retained span, so the package carries no omission record.  No retained line contains a citation, so the wrapper carries no bibliography and no bibliography entry.  No retained line contains a figure environment, an image-inclusion command or drawing code, so no image is delivered and the package declares no asset dependency.  Editorial formatting only, in addition: an AMS \texttt{amsart} preamble that restates the article's own declarations and macros used by the retained lines, namely the environments \texttt{lemma}, \texttt{proposition} on separate counters, together with the standard packages those lines need.  The retained environments print in this document as Lemma 1, Proposition 1 in order of appearance, whereas the article prints the same items with the numbering of its own full document.  The complete native source of the article as distributed in the arXiv e-print for this exact version is bundled unchanged as \texttt{sources/197148/source0.tex}.
\begin{lemma}[Fano via KL from the mean distribution]
\label{lem:fano-kl-mean}
Let $M\ge 2$ and let $V\sim\Unif([M])$.
For each $v\in[M]$ let $P_v$ be a distribution on $\Y$, and draw
$Y\mid(V=v)\sim P_v$.
Let $\widehat V=f(Y)$ be any estimator and set $P_e:=\Pr(\widehat V\neq V)$.
Define the mean (mixture) distribution
\[
\bar P \ :=\ \frac1M\sum_{v=1}^M P_v.
\]
Then
\begin{equation}
\log M - \frac1M\sum_{v=1}^M \KL(P_v\|\bar P)
\ \le\ h(P_e)+P_e\log(M-1),
\label{eq:fano-kl-mean-tight}
\end{equation}
where $h(p)=-p\log p-(1-p)\log(1-p)$ is the binary entropy.

In particular,
\begin{equation}
P_e
\ \ge\
1-\frac{\frac1M\sum_{v=1}^M \KL(P_v\|\bar P)+\log 2}{\log M}.
\label{eq:fano-kl-mean-simplified}
\end{equation}
\end{lemma}

\begin{lemma}[KL for biased product Bernoullis]
\label{lem:kl-product-bernoulli}
Fix $d\in\mathbb N$ and $\nu\in(0,\tfrac14)$.
For each $\theta\in\{\pm1\}^d$ define the product distribution
\[
P_\theta \;:=\; \bigotimes_{i=1}^d \Ber\!\left(\tfrac12+\theta_i\nu\right).
\]
Then for any $\theta,\theta'\in\{\pm1\}^d$,
\[
\KL(P_\theta\|P_{\theta'})
\;=\;
\Delta(\theta,\theta')\cdot D_\nu,
\]
where $\Delta(\theta,\theta'):=|\{i:\theta_i\neq\theta_i'\}|$ is the Hamming distance and
\[
D_\nu
\;:=\;
\KL\!\left(\Ber\!\left(\tfrac12+\nu\right)\Big\|\Ber\!\left(\tfrac12-\nu\right)\right)
=
2\nu\log\!\left(\frac{\tfrac12+\nu}{\tfrac12-\nu}\right)
=
2\nu\log\!\left(\frac{1+2\nu}{1-2\nu}\right).
\]
Moreover, for all $\nu\in(0,\tfrac14)$,
\[
8\nu^2 \;\le\; D_\nu \;\le\; \frac{32}{3}\,\nu^2.
\]
\end{lemma}

\begin{lemma}[Hellinger separation for biased product Bernoullis]
\label{lem:hellinger-product-bernoulli}
Let $k\ge 1$ and define product measures on $\{0,1\}^k$ by
\[
P_+ := \Ber\!\left(\tfrac12+\nu\right)^{\otimes k},
\qquad
P_- := \Ber\!\left(\tfrac12-\nu\right)^{\otimes k},
\]
where $0<\nu<1/2$. Let
\[
H^2(P,Q)\ :=\ 2\bigl(1-\rho(P,Q)\bigr),
\qquad
\rho(P,Q)\ :=\ \sum_{x}\sqrt{P(x)Q(x)}
\]
denote the squared Hellinger distance and the Bhattacharyya coefficient.
Then for every $0<\eps\le 1$, if
\[
\nu \;\ge\; \sqrt{\frac{\eps}{2k}},
\]
we have
\[
H^2(P_+,P_-) \;\ge\; \eps .
\]
\end{lemma}

\begin{proposition}[$\Omega(d)$ samples for TV/uniform estimation over the $d$--cube]
\label{prop:Omega-d-product}
Let $\X=\{0,1\}^d$ and let $\F=2^\X$.
Fix $\eps\in(0,1/10)$ and $\delta\in(0,1/3)$.
Consider the family of \emph{product} distributions
\[
\Pc := \{P_\theta : \theta\in\{\pm1\}^d\},
\qquad
P_\theta := \bigotimes_{i=1}^d \Ber\!\left(\tfrac12+\theta_i \nu\right),
\]
where
\[
\nu \;:=\; 4\,\sqrt{\frac{\eps}{d}}.
\]
Assume $d$ is large enough so that $\nu<1/4$.

Suppose an estimator $\widehat P$ based on $n$ i.i.d.\ samples satisfies
\[
\Pr_{S\sim P^{\otimes n}}
\Bigl[
\sup_{F\in\F}|\widehat P(F)-P(F)| \le \eps
\Bigr]\ \ge\ 1-\delta
\qquad\text{for all }P\in\Pc.
\]
Then necessarily
\[
n \ \ge\ c\,\frac{d}{\eps},
\]
for a universal constant $c>0$.
In particular, for constant $\eps$, one must have $n=\Omega(d)$.
\end{proposition}

\begin{proof}
\medskip
\noindent\textbf{Step 1: a large code with Hamming distance $\ge d/4$.}
Let $\Theta\subseteq\{\pm1\}^d$ be such that for all distinct
$\theta,\theta'\in\Theta$,
\[
\Delta(\theta,\theta'):=|\{i:\theta_i\neq\theta_i'\}|\ \ge\ d/4,
\]
and $|\Theta|\ge 2^{c_1 d}$ for a universal constant $c_1>0$.
Such a set exists by the Gilbert--Varshamov bound.
We restrict attention to the subfamily $\{P_\theta:\theta\in\Theta\}$.

\medskip
\noindent\textbf{Step 2: separation in total variation.}
Fix distinct $\theta,\theta'\in\Theta$ and let $k:=\Delta(\theta,\theta')\ge d/4$.
Let $I:=\{i\in[d]:\theta_i\neq\theta_i'\}$, so $|I|=k$.
Since the Bhattacharyya coefficient factors over product measures,
and the marginals coincide on $[d]\setminus I$, we have
\[
H^2(P_\theta,P_{\theta'}) \;=\; H^2\!\left(\Ber(\tfrac12+\nu)^{\otimes k},
\Ber(\tfrac12-\nu)^{\otimes k}\right).
\]
By Lemma~\ref{lem:hellinger-product-bernoulli} applied with $\eps' := 8\eps$,
it suffices that
\[
\nu \ \ge\ \sqrt{\frac{\eps'}{2k}}
\ =\ \sqrt{\frac{8\eps}{2k}}
\ =\ 2\sqrt{\frac{\eps}{k}}
\ \le\ 4\sqrt{\frac{\eps}{d}},
\]
where the last inequality uses $k\ge d/4$.
With our choice $\nu = 4\sqrt{\eps/d}$, we obtain
\[
H^2(P_\theta,P_{\theta'}) \ \ge\ 8\eps .
\]
Since $\TV(P,Q)\ge \tfrac12 H^2(P,Q)$, it follows that
\[
\TV(P_\theta,P_{\theta'}) \ \ge\ 4\eps
\qquad
\text{for all distinct }\theta,\theta'\in\Theta.
\]

\medskip
\noindent\textbf{Step 3: uniform estimation implies decoding.}
For distributions on a finite domain,
\[
\TV(P,Q)=\sup_{F\subseteq\X}|P(F)-Q(F)|.
\]
Thus, by assumption, with probability at least $1-\delta$,
\[
\TV(\widehat P,P_\theta)\ \le\ \eps.
\]
Define the decoder
\[
\widehat\theta
\in \arg\min_{\theta'\in\Theta} \TV(\widehat P,P_{\theta'}).
\]
On the event $\TV(\widehat P,P_\theta)\le\eps$, separation implies
$\widehat\theta=\theta$: indeed, for $\theta'\neq\theta$,
\[
\TV(\widehat P,P_{\theta'})
\ \ge\ \TV(P_\theta,P_{\theta'})-\TV(\widehat P,P_\theta)
\ \ge\ 4\eps-\eps
\ >\ \eps.
\]
Hence
\begin{equation}\label{eq:decode-success}
\Pr[\widehat\theta=\theta]\ \ge\ 1-\delta.
\end{equation}

\medskip
\noindent\textbf{Step 4: Fano via KL from the mean distribution.}
Let $V$ be uniform on $\Theta$, and let
$S\mid(V=\theta)\sim(P_\theta)^{\otimes n}$.
Let $\widehat V=\widehat\theta(S)$.
By~\eqref{eq:decode-success},
\[
P_e:=\Pr[\widehat V\neq V]\ \le\ \delta.
\]
Let
\[
\bar Q \ :=\ \frac1{|\Theta|}\sum_{\theta\in\Theta}(P_\theta)^{\otimes n}
\]
be the mean distribution of $S$.
Applying Lemma~\ref{lem:fano-kl-mean} gives
\begin{equation}\label{eq:fano-mean}
\log|\Theta|-\frac1{|\Theta|}\sum_{\theta\in\Theta}
\KL\!\left((P_\theta)^{\otimes n}\middle\|\bar Q\right)
\ \le\ h(P_e)+P_e\log(|\Theta|-1).
\end{equation}
Using $P_e\le\delta$, $h(P_e)\le\log2$, and $\log(|\Theta|-1)\le\log|\Theta|$,
we obtain
\[
\frac1{|\Theta|}\sum_{\theta\in\Theta}
\KL\!\left((P_\theta)^{\otimes n}\middle\|\bar Q\right)
\ \ge\ (1-\delta)\log|\Theta|-\log2.
\]

\medskip
\noindent\textbf{Step 5: upper bound the KL from the mean.}
By convexity of KL in its second argument,
\[
\KL\!\left((P_\theta)^{\otimes n}\middle\|\bar Q\right)
\le
\frac1{|\Theta|}\sum_{\theta'\in\Theta}
\KL\!\left((P_\theta)^{\otimes n}\middle\|(P_{\theta'})^{\otimes n}\right).
\]
Averaging over $\theta$ yields
\[
\frac1{|\Theta|}\sum_{\theta\in\Theta}
\KL\!\left((P_\theta)^{\otimes n}\middle\|\bar Q\right)
\le
\frac{n}{|\Theta|^2}\sum_{\theta,\theta'\in\Theta}
\KL(P_\theta\|P_{\theta'}).
\]
By Lemma~\ref{lem:kl-product-bernoulli},
\[
\KL(P_\theta\|P_{\theta'})
=\Delta(\theta,\theta')\,D_\nu
\le d\cdot D_\nu.
\]
Using the upper bound $D_\nu\le \frac{32}{3}\nu^2$, we get
\[
\frac1{|\Theta|}\sum_{\theta\in\Theta}
\KL\!\left((P_\theta)^{\otimes n}\middle\|\bar Q\right)
\ \le\ n\cdot C\,d\,\nu^2
\]
for a universal constant $C>0$.

\medskip
\noindent\textbf{Step 6: conclude the sample complexity lower bound.}
Combining the lower bound from Step~4 with the upper bound from Step~5 and using
$\log|\Theta|\ge c_1 d$, we obtain
\[
n\cdot C\,d\,\nu^2
\;\ge\;
(1-\delta)\log|\Theta|-\log 2
\;\ge\;
(1-\delta)c_1 d-\log 2 .
\]
Since $\delta\le 1/3$, the right-hand side is at least $c_2 d$ for a universal
constant $c_2>0$ (absorbing the additive $-\log 2$ term into the constant for
all sufficiently large $d$). Hence
\[
n\cdot C\,d\,\nu^2 \;\ge\; c_2 d,
\]
and therefore
\[
n
\;\ge\;
\frac{c_2}{C\,\nu^2}.
\]
Recalling that $\nu = 4\sqrt{\eps/d}$ (so that $\nu^2 = 16\,\eps/d$), this yields
\[
n
\;\ge\;
\frac{c_2}{C}\cdot \frac{d}{16\,\eps}
\;=\;
\frac{c_2}{16C}\,\frac{d}{\eps}.
\]
Setting $c:=c_2/(16C)$ completes the proof.
\end{proof}

\end{document}
